\section{Experimental design}
\subsection{Datasets}
\begin{table}[!htbp]
\centering\small
\caption{Benchmark datasets organized across three distinct inductive-bias categories with domain-purposeful column definitions.}
\begin{tabularx}{\linewidth}{@{}l l l >{\raggedright\arraybackslash}X@{}}
\toprule
Category & Dataset & Qubits / Size & Benchmarking rationale, columns, and target\\
\midrule
\multirow{2}{*}{\shortstack[l]{1. Quantum State\\\& Phase Classification}} 
& PennyLane TFIM & 4 qubits; $n=66$ & Transverse-Field Ising Model 1D open chain ($L=4$) ground states across transverse field $h \in [0, 3.0]$. Target: \texttt{magnetic\_phase} (0: Ferromagnetic, 1: Paramagnetic). Features: 4 physical Pauli expectation observables ($\langle Z_0 Z_1 \rangle, \langle Z_1 Z_2 \rangle, \langle Z_2 Z_3 \rangle, \langle X_{\text{mean}} \rangle$).\\
& Qiskit Phase of Matter & 4 qubits; $n=200$ & Bond-alternating XXX Heisenberg spin chain ground states (Qiskit ML \cite{bermejo2024quantum}). Target: \texttt{topological\_phase} (0: Trivial Dimerized, 1: Topological SPT). Features: 4 Heisenberg bond correlation and dimerization order observables ($\langle \vec{S}_0 \cdot \vec{S}_1 \rangle, \langle \vec{S}_1 \cdot \vec{S}_2 \rangle, \langle \vec{S}_2 \cdot \vec{S}_3 \rangle, O_{\text{dimer}}$).\\
\midrule
\shortstack[l]{2. High-Dimensional\\Quantum Kernels}
& Havl\'i\v{c}ek Ad-Hoc & 2 qubits; $n=200$ & Havl\'i\v{c}ek et al.\ (Nature 2019 \cite{havlicek2019supervised}) quantum feature map benchmark with separation gap 0.3. Target: \texttt{quantum\_kernel\_class} (0: Class A, 1: Class B). Features: 2 phase coordinate angles $x_1, x_2 \in [0, 2\pi]$ mapped via non-linear entangling $U_{\Phi}(\vec{x})$.\\
\midrule
\multirow{2}{*}{\shortstack[l]{3. Classical Small-Data\\\& Low-Sample Benchmarks}}
& Breast Cancer Wisconsin & 4 qubits; $n=200$ & Clinical diagnostic oncology dataset (UCI / Scikit-Learn). Target: \texttt{diagnosis} (0: Malignant, 1: Benign). Features: 4 primary clinical cell nucleus morphological biomarkers (\texttt{mean\_radius}, \texttt{mean\_texture}, \texttt{mean\_perimeter}, \texttt{mean\_concave\_points}).\\
& MNIST Digits (0 vs.\ 1) & 4 qubits; $n=200$ & Low-sample handwritten digit recognition. Target: \texttt{digit\_label} (0: Digit 0, 1: Digit 1). Features: 4 morphological spatial image features (\texttt{center\_core\_intensity}, \texttt{outer\_perimeter\_intensity}, \texttt{vertical\_spine\_intensity}, \texttt{horizontal\_cross\_intensity}).\\
\bottomrule
\end{tabularx}
\label{tab:datasets}
\end{table}

\subsection{Exploratory Data Analysis (EDA) Protocol}
Prior to model training, systematic Exploratory Data Analysis (EDA) is performed across all five benchmark datasets. The EDA pipeline explicitly inspects:
\begin{enumerate}
  \item \textbf{Column schemas and target definitions:} Feature column names, storage data types (\texttt{float32}), non-null counts, explicit target column identification (\texttt{int64}), and class label semantics.
  \item \textbf{Sample and class balance:} Class count symmetry ($N_0/N_1$) to prevent majority-class gradient bias.
  \item \textbf{Feature distributional moments:} Empirical mean, standard deviation, min/median/max, skewness, and kurtosis across all feature dimensions.
  \item \textbf{Geometric separability:} Normalized centroid distance $S = \|\mu_0 - \mu_1\|_2 / (\bar{\sigma}_0 + \bar{\sigma}_1)$ and Fisher discriminant ratio to establish baseline linear separability.
  \item \textbf{Inter-feature correlation:} Pairwise Pearson correlation matrices to quantify feature redundancy and collinearity.
  \item \textbf{Spectral variance decomposition:} Principal component analysis (PCA) variance spectra assessing manifold dimensionality.
\end{enumerate}
EDA figures and structured statistical schemas are saved to \texttt{results/figures/eda\_all\_datasets.png} and \texttt{results/\allowbreak eda\_summary.json}.

\subsection{Preprocessing and split protocol}
The standard scaler is fitted using only the training samples and subsequently applied to held-out test data. The classical MLPs receive standardized continuous features $z \in \mathbb{R}^d$. For the quantum circuit, each standardized feature $z_i$ is mapped to a rotation angle:
\begin{equation}
  x_i^{(q)} = \pi\,\operatorname{clip}\!\left(\frac{z_i}{2.5},-1,1\right).
\end{equation}
This bounds inputs strictly to $[-\pi,\pi]$, preventing angle overflow and ensuring uniform phase coverage across the Bloch sphere.

Ten random seeds (0--9) are evaluated. Within each seed, the MLP and every QNN configuration receive identical training and test splits. The seed dictates the train/test partition, the dataset generation for synthetic tasks, and reproducible initialization offsets. Because the exact same data realization and partition are used by all models at a given seed, each primary comparison is paired by seed.
